\section{Discussion and Limitations}

\subsection{What the result demonstrates}

The main conclusion is architectural rather than brand-specific. In this retrospective packet, a constrained schedule-selection layer has lower regret than a raw forecast schedule and a conservative heuristic under a common LP/oracle. V2+ combines a deterministic candidate library, prior-only family evidence, and V2 fallback. Similar ordering appears in mean and median summaries and four adjacent windows.

The study also demonstrates disciplined mixed evidence. TFT, Poland features, the corrected RF suite, the time-separated Decision Transformer suite, and the differentiable suite do not replace the fallback. The historical random forest remains an in-packet diagnostic because its training candidates are exact timestamp-shifted copies of evaluation candidates. The genuine temporal DT suite either abstains or harms. The differentiable suite supplies a narrower positive result: transformer correction outperforms a matched MLP in 28/36 comparisons and often improves the raw schedule, while remaining worse than V2+ in all 72 runs. Reporting both sides prevents an architecture signal from becoming a false controller claim.

For practitioners, the regret decomposition suggests that modeling effort should be targeted. If no useful candidate exists, selector complexity will not help. If a useful candidate exists but is missed, prior-context features and selective prediction are plausible targets. If strict is already near oracle, abstention is appropriate. This decomposition is more actionable than a single forecast score.

\subsection{Statistical limitations}

The main slice contains 90 profile-date rows but only 18 distinct dates. Five configured profiles share one market price process, and the dates are temporally ordered. The four windows are adjacent sensitivity replays within a limited 2026 period. The paired date-block summaries inserted in Section 5 treat each date as a cluster containing all five profiles; they are more appropriate than row-wise IID intervals but still cannot establish behavior outside the period.

The canonical random-forest aggregate uses three nominal seeds, but the selected path and mean regret are identical across seeds, producing zero aggregate standard deviation. The legacy numerical p-value has no inferential interpretation and is not evidence of stochastic replication; the corrected aggregate stores a null p-value. The paper relies on packet-level lineage and explicit caveats rather than statistical significance.

The corrective temporal suites use three seeds only as model-stability checks and adjacent windows as protocol sensitivity. Their runs are not independent market replications. The differentiable suite similarly has 90 profile rows per run but only 18 market dates, and model rows reuse the same declared windows. Counts such as 28/36 and 15/18 describe protocol stability; they are not binomial trials from independent markets.

The RF outcome is sparse and clustered: four profile-row switches on one delivery date drive the entire difference. Three are wins and one is a tie; 0 of 4 exceeds the configured 150 UAH tail-loss threshold. This is zero observed modeled tail losses in one correlated market episode, not a zero-risk estimate. Threshold sensitivity was performed on the same exact-mirror packet. A later period could remove or reverse the opportunity.

\subsection{Data and market limitations}

The primary performance evidence is DAM. IDM support is a read-model capability, and the paper does not claim IDM bidding performance. The market context is Ukrainian; transferring the numerical result to another bidding zone would require local price, publication, currency, rules, and battery-context validation. Neighboring-market features require licensing, timezone/DST, publication-time, and domain-shift checks before they can become official training covariates.

The panel uses five configured BESS profiles rather than measured commercial sites. It varies capacity, power, efficiency, and SOC settings but does not prove fleet-scale external validity. Measured telemetry, sensor quality, outages, maintenance states, and operator constraints would add uncertainty absent from this replay.

Source-readiness evidence is incomplete for the V13 research gate. A first-seen observation cannot be assumed to be an official publication timestamp, and a source preview is not a submission receipt. These distinctions limit the claims but prevent a more serious provenance error.

\subsection{Physical-model limitations}

The LP uses a power-energy battery model with efficiency, SOC, power constraints, and a throughput degradation proxy. Initial SOC resets for each profile-date and terminal SOC is unconstrained, so reported decision value is conditional one-day evaluator value rather than a repeatable chained multi-day profit estimate. The formulation has no explicit charge/discharge complementarity constraint, and the compact bundle does not preserve enough hourly solver output to audit whether numerically simultaneous flows occur. It also does not capture temperature-dependent efficiency, cell imbalance, nonlinear voltage behavior, calendar ageing, detailed state of health, or manufacturer-specific control limits. Literature on electro-thermal and degradation-aware scheduling shows that these factors can change value \cite{hesse2019ageing,maheshwari2020degradation,kumtepeli2020arbitrage}.

The oracle is also conditional on this simplified model. It is an oracle for the configured LP and realized price vector, not a universal upper bound on realizable asset profit. A richer digital twin could change both selected and oracle schedules. The current strength is comparative consistency: every method faces the same abstraction.

\subsection{Method limitations}

V2+ searches a finite, hand-designed candidate library. It can miss feasible schedules outside the included families. Its prior-only selector can also miss a better available family, as the oracle-gap audit shows. A full differentiable decision-focused system might learn forecast parameters that better support optimization \cite{mandi2024dfl,sang2022essdfl}, while a predict-then-bid formulation could model strategic market interaction \cite{yi2025predictbid}. Those are future systems, not alternate labels for V2+.

The HF value-aligned path is finite and cannot claim LP optimality. Its frozen 158.71 UAH diagnostic is threshold-selected on the same mirrored packet used to build its training representation; it is not independent of the co-primary V2+ packet. The separate 32-day read-model audit has no realized-regret outcome. The paper retains both as clearly separated engineering diagnostics only.

Decision Transformer remains exploratory. The old 460.30 UAH value comes from mirrored training rows. The corrective 36-run suite uses genuine return conditioning and temporal separation but still produces no gain over V2+. This is evidence against the tested objectives and dataset, not a general DT conclusion. The differentiable price-transformer result suggests that hourly sequence structure is useful, but the small number of distinct price paths, relaxed/strict optimization mismatch, and finite V2+ schedule library remain important confounders.

\subsection{Operational and ethical limitations}

No human-subject data are used, but operator-facing energy recommendations can still create harm if presented as authorized commands. The interface must preserve source labels, uncertainty, fallback reasons, and the no-execution boundary. Human oversight is not a decorative confirmation button; the system has been designed so that there is no executable market or hardware action to confirm.

Economic values are backtest evidence, not promises of profit. They exclude parts of a production revenue stack and operational cost structure, including settlement effects, imbalance exposure, taxes, fees, downtime, and detailed ageing. Users should not make investment decisions from the reported UAH values alone.

\subsection{Future work}

The nearest research step remains prospective validation of the frozen v1.2 transformer challenger under the unchanged V2+ fallback and oracle scorer. The next protocol should add new market dates rather than tune the stored windows, preserve the forecast-loss transformer as declared, and require stable improvement over both raw and V2+ after costs. Terminal-SOC equality or value, chained SOC, and explicit non-simultaneous charge/discharge sensitivity should be aligned between relaxed training and strict evaluation before interpreting daily value as repeatable arbitrage value.

A second direction is candidate-library expansion. Teacher schedules from a richer LP or MIP could provide diverse but feasible actions. Distillation could predict schedule vectors or family scores, followed by deterministic projection and repair. Comparison against the optimizer must remain on frozen windows.

A third direction is stronger uncertainty and source context. TFT quantiles and governed neighboring-market covariates could be admitted only after point-in-time coverage and robustness gates. Better source-publication evidence would support more precise readiness claims, while market-submission receipts would remain a separate requirement.

Finally, user-specific operational constraints should enter the optimizer or repair layer explicitly: critical-load reserves, protected hours, minimum backup SOC, maintenance windows, tariff rules, and asset preferences. UI labels are not a substitute for formal constraints.
